% 标准型 2x1+x2+x3=4, x1+2x2+x4=4 的六个基本解（全部两列子集都线性无关）：
% B={3,4}:(0,0)  B={1,4}:(2,0)  B={2,3}:(0,2)  B={1,2}:(4/3,4/3) 可行；B={1,3}:(4,0) 有 x3=-4、B={2,4}:(0,4) 有 x4=-4，不可行。
\begin{tikzpicture}[x=11mm,y=11mm,
  note/.style={font=\FigureNoteFont\fontsize{7.5}{9}\selectfont,text=FigureMuted},
  feas/.style={fill=FigureBlue,draw=FigureBlue},
  infeas/.style={fill=white,draw=FigureRed,line width=.9pt}]
  \fill[FigureBlueFill] (0,0)--(2,0)--(4/3,4/3)--(0,2)--cycle;
  \draw[FigureStroke!70,line width=.6pt] (-0.15,4.3)--(2.2,-0.4);
  \draw[FigureStroke!70,line width=.6pt] (-0.3,2.15)--(4.6,-0.3);
  \draw[->,FigureStroke,line width=.6pt] (-0.3,0)--(4.8,0) node[right]{$x_1$};
  \draw[->,FigureStroke,line width=.6pt] (0,-0.3)--(0,4.6) node[above]{$x_2$};
  \foreach \p in {(0,0),(2,0),(0,2),(4/3,4/3)}{\path[feas] \p circle (1.8pt);}
  \foreach \p in {(4,0),(0,4)}{\path[infeas] \p circle (2pt);}
  \node[note,anchor=north east] at (-0.02,-0.02) {$B=\{3,4\}$};
  \node[note,anchor=north] at (2.15,-0.08) {$B=\{1,4\}$};
  \node[note,anchor=east] at (-0.06,2) {$B=\{2,3\}$};
  \node[note,anchor=south west] at (1.4,1.38) {$B=\{1,2\}$};
  \node[note,anchor=north,text=FigureRed!80!FigureInk] at (4,-0.08) {$B=\{1,3\}$};
  \node[note,anchor=north,text=FigureRed!80!FigureInk] at (4,-0.4) {$x_3=-4$};
  \node[note,anchor=west,text=FigureRed!80!FigureInk] at (0.1,4.1) {$B=\{2,4\}$，$x_4=-4$};
  \node[note,anchor=west] at (1.1,2.4) {$2x_1+x_2=4$（$x_3=0$）};
  \node[note,anchor=west] at (2.6,0.95) {$x_1+2x_2=4$（$x_4=0$）};
\end{tikzpicture}
